\section{Introduction}\label{sec:Intro}
Finite element analysis is a prevalent methodology for the precise computation of electric machines; nevertheless, it entails substantial execution time. To expedite the design process, analytical methods are frequently employed, providing a faster development cycle. These methods can be categorized into three types: potential-theoretic, hybrid, and magnetic equivalent circuits. The hybrid approach leverages potential-theoretic analysis and integrates the resultant magnetic flux into magnetic equivalent circuits using ideal current sources, thereby decoupling rotor-stator interactions \cite{Huang2014, Kano2010}. Although hybrid methods present advantages, they are beyond the scope of this paper, which emphasises rotor-stator interactions through air-gap reluctances. Given the continuous variation in the overlapping areas between rotor and stator segments, precise computation of node-specific cross-sectional areas is imperative for determining the air-gap reluctances. An approach to address this challenge is outlined in \cref{sec:primfac}.

The determination of the winding distribution is pivotal for the modeling of electric machines, as it exerts an influence on their electromagnetic performance. Literature highlights methodologies such as the slot star and the Tingley diagram, which, particularly for concentrated windings, are subject to specific constraints that impede their practical implementation \cite{MüllerBerechnung, Binder2018, Rocha20}. In \cref{sec:WindingDist} of this study, an alternative approach is presented that enables a straightforward determination of the winding distribution and allows for the immediate identification of the $\alpha$-axis in the stator-fixed coordinate system, thereby facilitating the alignment between the stator and rotor.

The presentation utilizes axial flux machines with a yokeless design, which are occasionally transformed into linear machines through radial sectioning and subsequent unfolding. However, the methodologies are equally applicable to radial flux and linear machines.